\subsection{Accuracy of All Models (RQ2)}
\label{sec:res-acc}
\TabMain
Table~\ref{tab:main} gives the main results. On the unified test set the teacher reaches a macro-F1 of \val{teacher.test.macro} [\val{teacher.test.ci}] and an accuracy of \val{teacher.test.acc}. This is far below its development score (\val{teacher.dev.macro}), because the test set covers languages, sources, and a period the teacher has never seen. StaticKD reaches a macro-F1 of \val{statickd.test.macro} [\val{statickd.test.ci}] and an accuracy of \val{statickd.test.acc}, which is \val{statickd.test.pct_teacher} of the teacher macro-F1. On the development set StaticKD keeps \val{statickd.dev.pct_teacher} of the teacher macro-F1 (\val{statickd.dev.macro} against \val{teacher.dev.macro}), and its agreement with the teacher over the \val{data.heldout.strict.langs} held-out languages is \val{statickd.heldout.acc}. A single student without ensembling reaches \val{single.test_macro.mean} $\pm$ \val{single.test_macro.sd} on the test set and \val{single.dev_macro.mean} $\pm$ \val{single.dev_macro.sd} on the development set (10 seeds).

Among the small distilled transformers, MiniLM-L6 (\val{minilm_kd.test.macro}), which starts from a distillation of XLM-R-large, lies between StaticKD and the teacher. DistilmBERT (\val{distilmbert_kd.test.macro}) is on par with StaticKD while being much more expensive. The two baselines of similar cost, fastText + KD and TF-IDF + SVM + KD, fall below StaticKD. fastText trained only on the GPT-4o labels of four languages fails outside its training languages, with a test macro-F1 of \val{ft_gpt_subword.test.macro} against \val{ft_gpt_subword.dev.macro} on the development set.

\TabSig
Table~\ref{tab:sig} reports paired tests with $\Delta$ = StaticKD $-$ comparison model on the test set. Against the teacher, $\Delta$macro-F1 is \val{sig.test.teacher.delta} [\val{sig.test.teacher.ci}] ($p$ \val{sig.test.teacher.p}), and against MiniLM-L6 it is \val{sig.test.minilm_kd.delta} [\val{sig.test.minilm_kd.ci}] ($p$ \val{sig.test.minilm_kd.p}). The difference to DistilmBERT, \val{sig.test.distilmbert_kd.delta} [\val{sig.test.distilmbert_kd.ci}], is not significant ($p$ \val{sig.test.distilmbert_kd.p}). StaticKD is correct on more test documents than DistilmBERT (McNemar $p$ \val{sig.test.distilmbert_kd.pm}), but the McNemar test treats documents as independent and ignores the dependence between articles of the same site. StaticKD outperforms the quantized fastText + KD by \val{sig.test.ft_kd_subword_q.delta} [\val{sig.test.ft_kd_subword_q.ci}] ($p$ \val{sig.test.ft_kd_subword_q.p}) and TF-IDF + SVM by \val{sig.test.tfidf_svm_kd.delta} [\val{sig.test.tfidf_svm_kd.ci}] ($p$ \val{sig.test.tfidf_svm_kd.p}).

These two advantages differ in strength when examined per language. StaticKD has a higher macro-F1 than fastText + KD in \val{lang.static_better_macro.ft_kd_subword_q} of \val{lang.n} languages (Wilcoxon $p$ \val{lang.wilcoxon_macro.ft_kd_subword_q.p}) and a higher accuracy in \val{lang.static_better.ft_kd_subword_q} (Wilcoxon $p$ \val{lang.wilcoxon.ft_kd_subword_q.p}). Against TF-IDF + SVM it has a higher macro-F1 in \val{lang.static_better_macro.tfidf_svm_kd} of \val{lang.n} languages (Wilcoxon $p$ \val{lang.wilcoxon_macro.tfidf_svm_kd.p}), but a higher accuracy in only \val{lang.static_better.tfidf_svm_kd} and a lower accuracy in \val{lang.static_worse.tfidf_svm_kd} (Wilcoxon $p$ \val{lang.wilcoxon.tfidf_svm_kd.p}). The advantage over TF-IDF + SVM is thus consistent across languages on the main metric, while the per-language accuracies of the two models are close. Against DistilmBERT neither metric differs consistently across languages (Wilcoxon $p$ \val{lang.wilcoxon_macro.distilmbert_kd.p} for macro-F1 and \val{lang.wilcoxon.distilmbert_kd.p} for accuracy). StaticKD beats the teacher in only \val{lang.static_better.teacher} of \val{lang.n} languages (Wilcoxon $p$ \val{lang.wilcoxon.teacher.p}).

By label origin (Table~\ref{tab:groups}, Fig.~\ref{fig:groups}), StaticKD retains \val{grp.label_origin.publishersection.pct} of the teacher score on articles with JSON-LD section labels, \val{grp.label_origin.publishercategor.pct} on publisher categories, and \val{grp.label_origin.humanannotation.pct} on human annotations. By the distillation data of the language, it retains \val{grp.distil_tier.cukup.pct} in languages with sufficient data, \val{grp.distil_tier.sedikit.pct} in languages with little data, and \val{grp.distil_tier.tidakada.pct} in languages with no distillation data at all. On a subset with at most five articles per site and cell (\val{sens.n} articles), the model ranking is \val{sens.rank_same} relative to the full test set (StaticKD \val{sens.statickd}, teacher \val{sens.teacher}).

\TabGroups

Table~\ref{tab:sources} splits the test set by source. On MN-DS, the only source annotated directly with IPTC topics, StaticKD keeps \val{src.mnds.pct} of the teacher score (\val{src.mnds.statickd} against \val{src.mnds.teacher}), and the teacher itself reaches its lowest macro-F1 of all sources there. The lowest retention is on MasakhaNEWS (\val{src.masakhanews.pct}), whose human labels follow a non-IPTC scheme and most of whose African languages are absent from the distillation data. Bangla is also absent from the distillation data. On Bangla news, fastText + KD and TF-IDF + SVM collapse to \val{src.bangla_plos.ft_kd_subword_q} and \val{src.bangla_plos.tfidf_svm_kd}, while StaticKD reaches \val{src.bangla_plos.statickd}. This gap shows the cross-lingual transfer that baselines without pre-trained embeddings lack.

\TabSources

\begin{figure*}[!t]
\centering
\includegraphics[width=0.78\textwidth]{figures/fig_groups.pdf}
\caption{Test macro-F1 by label origin (left) and by the amount of distillation data in the language of the article (right).}
\label{fig:groups}
\end{figure*}

\subsection{Validity of the Editorial Labels}
\label{sec:res-validity}
Section and category labels record newsroom decisions, so the question is whether they agree with human labels in the same language.

All English cells of the test set are filled with human-labeled sources. As a comparison within the same language, a separate sample of \val{valid.en.silver.n} English CC-NEWS articles with JSON-LD section labels was therefore drawn from \val{valid.en.silver.sites} sites outside the test set, at most 50 per label. The teacher reaches an accuracy of \val{valid.en.silver.acc} on the section labels (macro-F1 \val{valid.en.silver.macro}) and \val{valid.en.human.acc} on the MN-DS and MasakhaNEWS human labels (macro-F1 \val{valid.en.human.macro}). Accuracy on section labels is higher for \val{valid.en.labels_higher} of the 17 labels (Wilcoxon $p$ \val{valid.en.wilcoxon.p}), and the ranking of label difficulty correlates between the two label types ($\rho=\val{valid.en.label_rho}$, $p$ \val{valid.en.label_rho.p}). The hard labels, \emph{society}, \emph{human interest}, and \emph{religion}, are hard under both.

French offers a second comparison inside the test set. The teacher reaches an accuracy of \val{valid.fr.rubrik.acc} on \val{valid.fr.rubrik.n} articles with section labels and \val{valid.fr.anotasi.acc} on \val{valid.fr.anotasi.n} human-labeled MasakhaNEWS articles.

A last check looks at each source category, restricted to languages covered by XLM-R pre-training. The median agreement of the teacher with the mapped label per category is \val{mapcheck.masakhanews.median} for MasakhaNEWS, \val{mapcheck.mnds.median} for MN-DS, \val{mapcheck.l3cube_indic.median} for L3Cube, and \val{mapcheck.ccnews_jsonld.median} for CC-NEWS sections. The lowest agreement occurs for source categories that are broader than their IPTC target. For example, only \val{mapcheck.masakha_politics} of MasakhaNEWS \emph{politics} articles are predicted as \emph{politics}, with \val{mapcheck.masakha_politics.alt} as the main alternative. Agreement is also low for labels whose boundaries are unclear even to human annotators, such as \emph{society} in MN-DS, in line with the low inter-annotator agreement reported for that label~\cite{kuzman2025llm}. Such mismatches reflect differences between the category scheme of a source and IPTC. Since every model is scored against the same labels, comparisons between models are unaffected, and absolute scores read as agreement with the editorial categorization.

\subsection{Ablation, Temperature, and Ensembling (RQ3)}
\TabAblation
\TabAblationGroups
Table~\ref{tab:ablation} gives the contribution of each component with five seeds per variant. Multilingual distillation data matter most. Distilling on EMMediaTopic alone (4 languages) lowers the test macro-F1 by \val{abl.emm_only.test_macro.delta} and the held-out agreement by \val{abl.emm_only.heldout_acc.delta}. Random initialization changes the test macro-F1 by \val{abl.random_init.test_macro.delta} ($p$ \val{abl.random_init.test_macro.p}), so the cross-lingually aligned embeddings account for a large share of the accuracy.

Freezing the embeddings and training only the linear head on top of potion lowers the development score by \val{abl.frozen.dev_macro.delta} ($p$ \val{abl.frozen.dev_macro.p}) and the held-out agreement by \val{abl.frozen.heldout_acc.delta} ($p$ \val{abl.frozen.heldout_acc.p}). On the test set the same variant raises macro-F1 by \val{abl.frozen.test_macro.delta} ($p$ \val{abl.frozen.test_macro.p}). Table~\ref{tab:ablation-groups} shows that this gain comes from languages with little distillation data (\val{ablgrp.frozen.sedikit.delta}) and with none (\val{ablgrp.frozen.tidakada.delta}). In data-rich languages the frozen embeddings are weaker (\val{ablgrp.frozen.cukup.delta}).

Hard teacher labels change the test macro-F1 by \val{abl.hard.test_macro.delta} ($p$ \val{abl.hard.test_macro.p}) and the development score by \val{abl.hard.dev_macro.delta} ($p$ \val{abl.hard.dev_macro.p}), so they do not differ from soft labels. Reducing the data to 130 thousand documents changes the test macro-F1 by \val{abl.data130k.test_macro.delta} ($p$ \val{abl.data130k.test_macro.p}), removing token dropout by \val{abl.no_tdrop.test_macro.delta} ($p$ \val{abl.no_tdrop.test_macro.p}), and an MLP head that cannot be folded by \val{abl.mlp.test_macro.delta} ($p$ \val{abl.mlp.test_macro.p}).

Fig.~\ref{fig:temperature} shows the effect of temperature. Compared with $\tau=1$, the common choice $\tau=2$~\cite{hinton2015distilling} lowers the held-out agreement by \val{abl.tau2.heldout_acc.delta} ($p$ \val{abl.tau2.heldout_acc.p}). Its changes on the test macro-F1 (\val{abl.tau2.test_macro.delta}, $p$ \val{abl.tau2.test_macro.p}) and on the development set (\val{abl.tau2.dev_macro.delta}, $p$ \val{abl.tau2.dev_macro.p}) are not significant. With $\tau=4$ all three metrics drop, and on the test set the drop is \val{abl.tau4.test_macro.delta} ($p$ \val{abl.tau4.test_macro.p}). None of the four exploratory variants improves the test set: bigram hashing \val{abl.bigram.test_macro.delta}, idf pooling \val{abl.idf.test_macro.delta}, confidence weights \val{abl.conf.test_macro.delta}, and cross-entropy on GPT-4o labels \val{abl.gpt_ce.test_macro.delta}.

Growing the ensemble from one to three seeds moves the test macro-F1 from \val{ens.k1.test_macro} to \val{ens.k3.test_macro} and the held-out agreement from \val{ens.k1.heldout_acc} to \val{ens.k3.heldout_acc} (Fig.~\ref{fig:ensemble}). Ten seeds give \val{ens.k10.test_macro}. The gain is small, but it costs nothing at inference time because every ensemble size is still one table.

\begin{figure*}[!t]
\centering
\includegraphics[width=0.9\textwidth]{figures/fig_temperature.pdf}
\caption{Effect of the distillation temperature. Each dot is one seed; the line connects the means.}
\label{fig:temperature}
\end{figure*}

\begin{figure}[!t]
\centering
\includegraphics[width=0.85\columnwidth]{figures/fig_ensemble_size.pdf}
\caption{Effect of the table ensemble size (mean $\pm$ standard deviation over seed combinations).}
\label{fig:ensemble}
\end{figure}

\subsection{CPU Efficiency (RQ1)}
\label{sec:res-eff}
\TabEff
Table~\ref{tab:eff} and Figs.~\ref{fig:tradeoff} and~\ref{fig:latency} show the efficiency of all models. In PyTorch fp32 the teacher needs \val{eff.teacher.latency_ms_p50_1c}~ms per article on one core, with a throughput of \val{eff.teacher.docs_per_sec_1c} articles per second. Its best CPU configuration, ONNX int8, needs \val{eff.teacher_onnx_int8.latency_ms_p50_1c}~ms with a throughput of \val{eff.teacher_onnx_int8.docs_per_sec_1c} articles per second. Classifying one million articles per day would take \val{eff.teacher_onnx_int8.cores1M} cores (ONNX int8) to \val{eff.teacher.cores1M} cores (PyTorch) for this step alone.

StaticKD needs \val{eff.statickd.latency_ms_p50_1c}~ms per article and processes \val{eff.statickd.docs_per_sec_1c} articles per second per core. This is about \val{eff.statickd.speedup_onnx} times faster than the ONNX int8 teacher, \val{eff.statickd.speedup_pt} times faster than the PyTorch teacher, and \val{eff.speed_vs_minilm} times faster than MiniLM-L6. Its file is \val{eff.size_ratio_onnx} times smaller than the ONNX int8 teacher, and one million articles per day need \val{eff.statickd.cores1M} cores. Its model RAM is \val{eff.statickd.ram}~MB, mostly for the tokenizer structure. The compact variant reduces it to \val{eff.statickd_compact.ram}~MB. With one thread, tokenization (\val{prof.statickd.tokens} tokens per article on average) takes \val{prof.statickd.tokshare} of the prediction time, while table pooling alone runs at \val{prof.statickd.pool_dps} articles per second. The speed limit of StaticKD is thus set by the tokenizer implementation.

StaticKD is not the fastest model in the table. The quantized fastText + KD is slightly faster (\val{eff.ft_kd_subword_q.docs_per_sec_1c} articles per second), smaller (\val{eff.ft_kd_subword_q.size}~MB), and uses less RAM (\val{eff.ft_kd_subword_q.ram}~MB), but its test macro-F1 is \val{ft_kd_subword_q.test.macro} against \val{statickd.test.macro} for StaticKD. Among the models with comparable latency, namely StaticKD, its compact variant, fastText, and TF-IDF + SVM, StaticKD is the most accurate (Fig.~\ref{fig:tradeoff}).

\begin{figure}[!t]
\centering
\includegraphics[width=\columnwidth]{figures/fig_tradeoff.pdf}
\caption{Test macro-F1 on the unified test set against single-core CPU throughput (log scale).}
\label{fig:tradeoff}
\end{figure}

\begin{figure*}[!t]
\centering
\includegraphics[width=0.9\textwidth]{figures/fig_latency_ram_size.pdf}
\caption{Latency per article, model RAM, and file size (log scale).}
\label{fig:latency}
\end{figure*}

\subsection{Analysis by Language and Label (RQ4)}
\TabFactors
The median per-language accuracy gap between StaticKD and the teacher is \val{lang.gap.median} (Fig.~\ref{fig:lang}). On the held-out set (\val{fac.langs} languages with at least 20 documents), the agreement of StaticKD correlates with the log amount of distillation data ($\rho=\val{fac.data.rho}$, $p$ \val{fac.data.p}) and with teacher confidence ($\rho=\val{fac.conf.rho}$, $p$ \val{fac.conf.p}), but does not differ across scripts (Kruskal--Wallis $H=\val{fac.kw.H}$, $p$ \val{fac.kw.p}; Table~\ref{tab:factors}). On languages outside XLM-R pre-training that are also absent from the distillation data (\val{lang.outxlmr.langs}), StaticKD reaches a mean accuracy of only \val{lang.outxlmr.static_acc}, while the teacher reaches \val{lang.outxlmr.teacher_acc}.

In \val{lang.vlowcov.n} languages (\val{lang.vlowcov.langs}), fewer than half of the token occurrences belong to tokens seen at least 10 times in the distillation data (at most \val{lang.vlowcov.cov_max}). StaticKD still reaches a mean accuracy of \val{lang.vlowcov.static_acc} in these languages (teacher \val{lang.vlowcov.teacher_acc}). The compact variant removes those tokens from the tokenizer, and its accuracy in these languages falls to \val{lang.vlowcov.compact_acc}. In the remaining languages the compact variant reaches \val{lang.restcov.compact_acc} against \val{lang.restcov.static_acc} for the full StaticKD. Section~\ref{sec:discussion} discusses the mechanism behind this pattern.

\begin{figure*}[!t]
\centering
\includegraphics[width=0.9\textwidth]{figures/fig_per_language.pdf}
\caption{Left: StaticKD accuracy against teacher accuracy per test language (color shows the distillation data of the language, size the number of articles). Right: held-out agreement per language against the amount of distillation data for StaticKD and fastText + KD.}
\label{fig:lang}
\end{figure*}

Per label (Fig.~\ref{fig:label}), the hardest labels for StaticKD are \val{label.hardest}, and the easiest are \val{label.easiest}. The largest gaps to the teacher occur for \val{label.biggest_gap}, and the most frequent confusions are \val{label.top_confusions}.

\begin{figure*}[!t]
\centering
\includegraphics[width=0.9\textwidth]{figures/fig_per_label.pdf}
\caption{F1 per label on the test set (left) and the StaticKD confusion matrix normalized per reference label (right).}
\label{fig:label}
\end{figure*}

\subsection{Calibration and Cascade Mode}
\TabCascade
The ECE of StaticKD is \val{ece.statickd.dev} on the development set and \val{ece.statickd.test} on the test set, lower than that of the teacher (\val{ece.teacher.dev} and \val{ece.teacher.test}). Both models become overconfident on the test set, which lies outside the training distribution (Fig.~\ref{fig:calib}). In cascade mode, StaticKD decides the articles on which its confidence is at least $\theta$ and sends the rest to the teacher (Table~\ref{tab:cascade}). With $\theta=0.8$, StaticKD decides \val{casc.08.coverage} of the articles with an accuracy of \val{casc.08.acc_covered}. The cascade reaches an overall accuracy of \val{casc.08.accuracy} and a macro-F1 of \val{casc.08.macro_f1}, against \val{casc.101.accuracy} and \val{casc.101.macro_f1} for the teacher alone, at an effective throughput of \val{casc.08.docs_per_sec} articles per second per core against \val{casc.101.docs_per_sec}. Sending about two fifths of the articles to the teacher is enough for the cascade to match the teacher accuracy.

\begin{figure}[!t]
\centering
\includegraphics[width=0.8\columnwidth]{figures/fig_calibration.pdf}
\caption{Reliability diagram of StaticKD on the development and test sets.}
\label{fig:calib}
\end{figure}
